% Generated by scripts/assemble.py; corrections belong in source/corrections.json.
\examhead{189}{陕西师范大学 2026 年数学分析试卷}{178}
% p179-b14; source PDF page 179
\par\Needspace{5\baselineskip}
\blockanchor{p179-b14}
\majorquestion{一}计算题. 每题6分, 共30分.\par
% p179-b15; source PDF page 179
\blockanchor{p179-b15}
\question{1}求函数  $\frac{1}{x^{2}+3x+2}$  在 x=-4 处的幂级数展开式和收敛域.\par
% p179-b16; source PDF page 179
\blockanchor{p179-b16}
\question{2}设函数 $y = y(x)$ 由参数方程 $\left\{ \begin{array}{l} x = 2e^{t} + t + 1 \\ y = 4(t - 1)e^{t} + t^{2} \end{array} \right.$ 确定, 计算 $\frac{\mathrm{d}^2y}{\mathrm{d}x^2}$.\par
% p179-b17; source PDF page 179
\blockanchor{p179-b17}
\question{3}设 $f(x)$ 的一个原函数为 $F(x) > 0$, 且 $F(0) = 1$, 当 $x \geq 0$ 时, 有 $f(x)F(x) = \sin^22x$, 求 $f(x)$.\par
% p179-b18; source PDF page 179
\par\Needspace{6\baselineskip}
\blockanchor{p179-b18}
\question{4}计算曲线积分\par
\nopagebreak[4]
% p179-b19; source PDF page 179
\blockanchor{p179-b19}
\displayformula{\int_ {L} [ 2 (x ^ {2} - x) e ^ {2 x} - x y ] \mathrm{d} x + [ x ^ {2} - (y + 2) e ^ {y} ] \mathrm{d} y.}
% p179-b20; source PDF page 179
\blockanchor{p179-b20}
式中 L 是从  $(0,0)$  经曲线  $y = x^{2} - 2x$  到点 A(4,8) 的一段弧.\par
% p179-b21; source PDF page 179
\par\Needspace{6\baselineskip}
\blockanchor{p179-b21}
\question{5}计算曲面积分\par
\nopagebreak[4]
% p179-b22; source PDF page 179
\blockanchor{p179-b22}
\displayformula{\iint_ {\Sigma} x \mathrm{d} y \mathrm{d} z + y \mathrm{d} z \mathrm{d} x + z \mathrm{d} x \mathrm{d} y.}
% p179-b23; source PDF page 179
\blockanchor{p179-b23}
其中 $\Sigma$ 为上半球面 $z = \sqrt{R^2 - x^2 - y^2}$ ($R > 0$) 的上侧.\par
% p179-b24; source PDF page 179
\par\Needspace{5\baselineskip}
\blockanchor{p179-b24}
\majorquestion{二}解答题. 每题 10 分, 共 30 分.\par
% p179-b25; source PDF page 179
\blockanchor{p179-b25}
\question{6}解答如下问题:\par
% p179-b26; source PDF page 179
\blockanchor{p179-b26}
(1) 讨论函数 $f(x) = (x - 1)(x - 2)D(x)$ 的连续点和间断点, 并判断间断点的类型, 其中 $D(x)$ 是狄利克雷函数\par
% p179-b27; source PDF page 179
\blockanchor{p179-b27}
$D(x)=\left\{\begin{aligned}&1,&x \text{为有理数};\\ &0,&x \text{为无理数}.\end{aligned}\right.$\par
% p180-b00; source PDF page 180
\blockanchor{p180-b00}
(2) 给出实数集 $\mathbb{R}$ 上只有 3 个连续点的函数.\par
% p180-b01; source PDF page 180
\blockanchor{p180-b01}
\question{7}利用交换积分次序计算 \displayformula{\int_{1}^{2} \mathrm{d}x \int_{\sqrt{x}}^{x} \sin \frac{\pi x}{2y} \mathrm{d}y + \int_{2}^{4} \mathrm{d}x \int_{\sqrt{x}}^{2} \sin \frac{\pi x}{2y} \mathrm{d}y.}\par
% p180-b02; source PDF page 180
\blockanchor{p180-b02}
\question{8}讨论反常积分 $I = \int_{0}^{\frac{\pi}{2}} \ln \cos x \, \mathrm{d}x$ 的敛散性, 若收敛, 计算 $I$ 的值.\par
% p180-b03; source PDF page 180
\par\Needspace{5\baselineskip}
\blockanchor{p180-b03}
\majorquestion{三}证明题. 每题 15 分, 共 90 分.\par
% p180-b04; source PDF page 180
\blockanchor{p180-b04}
\question{9}设 $f(x)$ 在 $[0,1]$ 上连续, 在 $(0,1)$ 内可导, 且 $\int_0^1 f(x)\mathrm{d}x = 0$, 记 $F(x) = \int_0^x xf(t)\mathrm{d}t$.\par
% p180-b05; source PDF page 180
\blockanchor{p180-b05}
(1) 求  $F'(x)$ .\par
% p180-b06; source PDF page 180
\blockanchor{p180-b06}
(2) 证明: 在 $(0,1)$ 内至少存在一点 $\xi$, 使得 $\int_{0}^{\xi} f(x) \mathrm{d}x = -\xi f(\xi)$.\par
% p180-b07; source PDF page 180
\blockanchor{p180-b07}
(3) 证明: 在 $(0,1)$ 内至少存在一点 $x_0$, 使得 $2f(x_0) + x_0 f'(x_0) = 0$.\par
% p180-b08; source PDF page 180
\blockanchor{p180-b08}
\question{10}设 $a_{n} > 0, b_{n} > 0$, 证明:\par
% p180-b09; source PDF page 180
\blockanchor{p180-b09}
(1) 若存在  $\alpha > 0$ ，使得  $\frac{b_{n}}{b_{n+1}}a_{n} - a_{n+1} \geq \alpha (n = 1, 2, \cdots)$ ，则级数  $\sum_{n=1}^{\infty} b_{n}$  收敛.\par
% p180-b10; source PDF page 180
\blockanchor{p180-b10}
(2) 若  $\sum_{n=1}^{\infty}\frac{1}{a_{n}}$  发散, 且  $\frac{b_{n}}{b_{n+1}}a_{n}-a_{n+1}\leq0(n=1,2,\cdots)$ , 则级数  $\sum_{n=1}^{\infty}b_{n}$  发散.\par
% p180-b11; source PDF page 180
\blockanchor{p180-b11}
\question{11}设 $x \geq 0, n$ 为正整数, $f(x) = \int_{0}^{x} (t - t^2) \sin^{2n} t \, \mathrm{d}t$, 证明: $f(x) \leq \frac{1}{(2n + 2)(2n + 3)}$.\par
% p180-b12; source PDF page 180
\par\Needspace{6\baselineskip}
\blockanchor{p180-b12}
\question{12}设 $f(x)$ 是 $[0,1]$ 上的正值连续函数, 且 $f(x)$ 单调递减, 证明:\par
\nopagebreak[4]
% p180-b13; source PDF page 180
\blockanchor{p180-b13}
\displayformula{\frac {\int_ {0} ^ {1} x f ^ {2} (x) \mathrm{d} x}{\int_ {0} ^ {1} x f (x) \mathrm{d} x} \leq \frac {\int_ {0} ^ {1} f ^ {2} (x) \mathrm{d} x}{\int_ {0} ^ {1} f (x) \mathrm{d} x}.}
% p180-b14; source PDF page 180
\blockanchor{p180-b14}
\question{13}设 $S(x) = \sum_{n=1}^{\infty} ne^{-nx}$, 其中 $x > 0$, 求证:\par
% p180-b15; source PDF page 180
\blockanchor{p180-b15}
(1) 函数项级数 $\sum_{n=1}^{\infty} ne^{-nx}$ 在 $(0, +\infty)$ 上不一致收敛.\par
% p180-b16; source PDF page 180
\blockanchor{p180-b16}
(2) 和函数 $S(x)$ 在 $(0, +\infty)$ 上连续可导.\par
% p180-b17; source PDF page 180
\blockanchor{p180-b17}
\question{14}证明: 若在区间 $I$ 上函数列 $\{f_n(x)\}$ 一致收敛, 且对每个正整数 $n$, $f_n(x)$ 都在区间 $I$ 上有界, 则函数列 $\{f_n(x)\}$ 在 $I$ 上一致有界.\par
